\documentclass[12pt,a4paper,oneside]{article}
\usepackage[brazil]{babel}
\usepackage[utf8]{inputenc}
\usepackage{olymp}

\setcounter{problem}{0}

\begin{document}

\begin{problem}{Maior lote da rua}{rua}{2}

Talvez você já tenha notado que, na maioria das ruas, a numeração das casas varia de 12 em 12.
Isso acontece por dois motivos:

\begin{itemize}
\item no Brasil, a numeração das casas segue a distância em metros desde o início da rua;
assim, se uma casa tem 12 metros de frente e número 300, a casa seguinte terá número 312.
\item por várias décadas, os lotes na maioria das cidades do Brasil foram criados com 12 m de frente.
\end{itemize}

Entretanto, há casas construídas em meio lote e há outras em lotes maiores, então a numeração não é sempre de 12 em 12, já que o tamanho do lote pode variar.

Faça um programa para calcular o tamanho do maior lote em uma rua, considerando como tamanho do lote a diferença entre números de casas consecutivas.

%\Notes
%
%Observações, se tiver.

\InputFile

A entrada começa com uma linha contendo um número inteiro $T$, indicando o número de casas na rua.
A linha seguinte contém $T$ números inteiros $N_i$, em ordem crescente, indicado o número de cada casa.
Restrições: $2 \leq T \leq 1000$, $1 \leq N_i \leq 100000$.

\OutputFile

Escreva o tamanho do maior lote da rua.

%\Notes

\Example

\begin{example}
\exmp{
5
20 32 44 68 70
}{
24
}%
\end{example}

\begin{example}
\exmp{
6
10 20 50 60 70 80
}{
30
}%
\end{example}

\begin{example}
\exmp{
2
100 115
}{
15
}%
\end{example}

%Comentários sobre o exemplo, se necessário.

\end{problem}

\end{document}